\documentclass{standalone}
%\usepackage{fullpage}
\usepackage{tikz}
\usepackage{graphicx}

\usetikzlibrary{positioning}
%\usetikzlibrary{shapes.multipart}

\tikzstyle{rootClass}=[rectangle,black,draw, fill=orange!30, text width=5 cm]
\tikzstyle{essentialClass}=[rectangle,black,draw, fill=green!30, text width=5 cm]
\tikzstyle{exampleClass}=[rectangle,black,draw, fill=red!30, text width=5 cm]
\tikzstyle{narrowBox}=[text width=5 cm]
\tikzstyle{wideBox}=[text width=13 cm]

\tikzstyle{inheritedFrom}=[->,black, line width=2]
\tikzstyle{instantiatedIn}=[-,black, line width=2] 
%Note that A instantiatedIn B relationship reduces flexibility as the 
%children of A cannot be known to B now.

\newcommand{\nodeTextWidthRatio}{0.95}
\newcommand{\spaceafter}{\vspace{0mm}}
\newcommand{\method}[1]{\spaceafter\textbf{\texttt{\detokenize{#1}}}\\}
\newcommand{\desc}[1]{\ \ \ \texttt{\detokenize{#1}}\\}
\newcommand{\className}[1]{\large \textbf{\detokenize{#1}}}
\newcommand{\data}[1]{\spaceafter\textbf{\texttt{\detokenize{#1}}}\\}
\newcommand{\code}[1]{\verb+ #1 +\\}

\def\showExampleClasses{1}
\def\showEssentialClasses{1}
%Note that you can't have showExampleClasses=1 with showEssentialClasses=0
%as the former nodes refer to the latter and this will cause error.

\begin{document}

\begin{tikzpicture}


%NOTE: you can't refer to exampleClass and essentialClass nodes
% in definition of roots. Order of definition matters for that and
%other type of nodes come later. 
\node[rootClass, wideBox] (World) {
	\begin{tabular}{p{\nodeTextWidthRatio\textwidth}} 

\className{World}
\\  \hline
\data{objects: dict[ObjectId, Object]}
\data{__vis_data : dict}
\data{__creation_ts: float}
\data{__num_evolutions: int}
\data{__current_time_ms: int}
\data{__vis_output_filename: str}
\data{__duration_sec: int}
\data{__vis_frame_interval_ms: float}
\data{__next_frame_time_msec: float}
 \\ \hline 
\method{constructor(map: Map, vis_filename) }

\method{evolve(delta-t)}
\desc{calls evolve() of all objects which have no owner, takes care of offspring objects, and finally kills objects which ask for it}

\method{run()}
\desc{manages delta-t (how small it should be), manages intersections using intersect(), and calls evolve()}

\method{intersect() -> tuple}
\desc{instantiates InIn for each pair of objects and calls InIn.intersect() to evaluate the intersection}

\method{add_object(new_objects: dict[ObjectId, Object]) -> None}

\method{delete_expired_objects() -> None}

\method{register_intersections(intersection_result: InInType) -> None}

\method{pick_delta_t() -> float}

\method{update_visualization_json()}

\method{dump_all_shapes_info()}

\method{dump_vis_data_to_file()}

\method{visualize()}
\end{tabular} 


};
\node[rootClass, wideBox] (Object) [below=of World] {
	\begin{tabular}{p{\nodeTextWidthRatio\textwidth}} 
\className{Object}
\\ \hline
\data{oid: objectID}
\desc{every object needs an ID to be tracable in map and in state}
\data{name: string}
\desc{merely a name, irrelevant to the simulation}
\data{shape: Shape}
\desc{which could be empty if the object is an owner}
\data{position: Position}
\desc{position and orientation of an anchor point of the object}
\data{__previous_position: Position}
\desc{to keep the previous state so we can revert}
\data{owner_object: ObjectID}
\data{dependent_objects: dict[ObjectId, Object]}
\data{_latest_intersections: list[IntersectionInstance]}
\data{_infinitesimal_intersection_occured: bool}
\\ \hline
\method{evolve(delta-t) -> list[Object]: offspring-objects}
\desc{changes the state (position, internal attributes, etc) of the object}
\desc{trvial evolution: when the object never changes state}
\desc{offsprings are the possibly non-physical objects required to accomplish something.}

\method{bounding-box() -> Box}
\desc{returns a box which contains the whole object. used to optimize intersection evaluation}

\method{get-required-delta-t()}
\desc{calculates the delta-t it requires to operate. By default 0, meaning no requirement}

\method{time-to-die() -> bool}
\desc{tells the World if it wants to be eliminated. This might be where Agent Smith cheated the matrix!}

\method{add_dependent_object(obj: Object)}

\method{set_intersections(intersections: list[in_in.IntersectionInstance]) -> None}

\method{update_position(new_position)-> None}

\method{revert_position() -> None}

\method{visualize()}
\desc{returns the information required for visualization}

\method{dump_shape_info() - >None}

\method{is_evolvable() -> bool}
\end{tabular} 

};

\node[rootClass, wideBox ] (Map) [right=of World] {
	\begin{tabular}{p{\nodeTextWidthRatio\textwidth}} 
\className{Map}
\\  \hline
\data{next_available_id}
 \\ \hline 

 \method{parse_map(self, filename: str) -> dict[ObjectId, Object]}

 \method{get_next_id(self) -> ObjectId}

 \method{instantiate_object(obj_json, new_id: ObjectId, name: string, owner: Object) -> Object}

\method{get_shape(obj_json) -> Shape}

\method{get_position(obj_json) -> Shape}

\method{get_objects(self, obj_map: dict[ObjectId, Object], parsed, owner: Object) -> dict[ObjectId, Object]}
\end{tabular} 


};

\node[rootClass ] (Shape) [right=of Map] {
	\begin{tabular}{p{\nodeTextWidthRatio\textwidth}} 
\className{Shape}
\\ \hline
\\ \hline
\method{bounding_box()}
\desc{returns a x-y plane bounding box. Can be done using a generalized algorithm, no implemented only in the parent class.}

\method{dump_info()}
\end{tabular} 


};

\node[rootClass ] (Position) [above=of Shape] {
	\begin{tabular}{p{\nodeTextWidthRatio\textwidth}} 
\className{Position}

\desc{position and orientation}
\\  \hline
\data{x,y,z}

\data{phi,theta}
 \\ \hline 
\end{tabular} 


};

\node[rootClass,wideBox ] (InIn) [left=of World] {
	\begin{tabular}{p{\nodeTextWidthRatio\textwidth}} 
\className{IntersectionInstance}
\\ \hline
\data{object1}

\data{object2}

\data{__does_intersect = False}

\data{__is_infinitesimal = False }

\data{_intersection_points = [] }
\\ \hline
\method{intersect() -> None}

\desc{the entry point for performing the intersection algorithm on its two
objects. 
No return value, the other methods do the job.
The results are then taken by calling methods such as does_intersect,
is_infinitesimal and so on.}

\method{does_intersect() -> bool}

\method{cylinder_cube()}

\method{line_circle() -> list, float}

\method{rectangle_lines() -> list}

\method{rotation_matrix_2d()}

\method{horizontal_line_segment_intersection()}

\method{cube_cube()}

\method{cylinder_cylinder()}

\method{is_infinitesimal()}

\method{get_intersection_point()}
\end{tabular}
};

% \node[rootClass ] (WorldState) [right=of World] {
% 	\begin{tabular}{p{\nodeTextWidthRatio\textwidth}} 
\className{WorldState}

\desc{use: e.g. debug}

\desc{the simulator works in delta-t steps. A worldState can be recorded either before or after evolution of all the objects. A partial worldState is invalid unless every object state is flagged as completed or not. The reason is that the world cannot be initialized with a snapshot of a partially completed worldState without knowing which objects have already evolved in that snapshot.}
\\  \hline
\desc{A data structure which keeps a set of states for every objectID}
 \\ \hline 
\method{load(file)}

\method{dump(file)}

\method{extractState(objectID)}

\method{recordState(objects[])}

\desc{    For object in objects:}

\desc{         Get and record object.dumpState()}
\end{tabular} 


% };

%\draw [instantiatedIn] (Object) -- (World);

\if\showEssentialClasses1
	%NOTE: you can't refer to exampleClass nodes in definition of 
	%essentials or roots. Order of definition matters for that and
	%exampleClasses come later. 
	% \node[essentialClass ] (XMLmap) [below=of Map] {
	% 	\begin{tabular}{p{\nodeTextWidthRatio\textwidth}} 
\className{XMLmap}
\\  \hline
 \\ \hline 
\end{tabular} 


	% };

	% \node[essentialClass ] (DiscretizedShape) [below=of Shape] {
	% 	\begin{tabular}{p{\nodeTextWidthRatio\textwidth}} 
\className{DiscretizedShape}

\desc{takes a mesh as input. Mesh given in a file or something? }
\\ \hline
\\ \hline
\end{tabular} 


	% };

	% \node[essentialClass ] (DiscretizedShape2D) [below=of DiscretizedShape] {
	% 	\begin{tabular}{p{\nodeTextWidthRatio\textwidth}} 
\className{DiscretizedShape2D}

\desc{takes bitmap and if required extends it to a mesh of infinitely long pieces along z-axis}
\\ \hline
\\ \hline
\end{tabular} 


	% };

	\node[essentialClass, wideBox ] (RigidPhysicalObject) [right=of Object] {
		\input{nodes/essentialClasses/RigidPhysicalObject}
	};

	\node[essentialClass ] (Robot) [left=of Object] {
		\begin{tabular}{p{\nodeTextWidthRatio\textwidth}} 
\className{Robot}
\\ \hline
\\ \hline
\desc{evolve() is the heart of Robot actions.}
\end{tabular} 


	};

	\node[essentialClass ] (Sensor) [left=of Robot, below=of InIn] {
		\begin{tabular}{p{\nodeTextWidthRatio\textwidth}} 
\className{Sensor}

\desc{inherently an Object, so needs shape, position, evolve(), etc.}
\\ \hline
\\ \hline
\method{sense()} 
\end{tabular} 


	};

	\draw [inheritedFrom] (Robot) -- (Object);
	\draw [inheritedFrom] (Sensor) -- (Object);
	\draw [inheritedFrom] (RigidPhysicalObject) -- (Object);
	% \draw [inheritedFrom] (DiscretizedShape) -- (Shape);
	% \draw [inheritedFrom] (DiscretizedShape2D) -- (DiscretizedShape);
	% \draw [inheritedFrom] (XMLmap) -- (Map);
\fi

\if\showExampleClasses1
	% \node[exampleClass ] (Disk) [below=of Shape, right=of DiscretizedShape] {
	% 	\begin{tabular}{p{\nodeTextWidthRatio\textwidth}} 
\className{Disk}
\\ \hline
\data{axis1}

\data{axis2}

\data{Coordinate center}
\\ \hline
\end{tabular} 


	% };

	\node[exampleClass ] (Cylinder) [below=of Shape] {
		\begin{tabular}{p{\nodeTextWidthRatio\textwidth}} 
\className{Cylinder}
\\ \hline
\\ \hline
\end{tabular} 


	};

	\node[exampleClass ] (Cube) [right=of Cylinder] {
		\begin{tabular}{p{\nodeTextWidthRatio\textwidth}} 
\className{Cube}
\\ \hline
\\ \hline
\end{tabular} 


	};
	% \node[exampleClass ] (Wheel) [below=of Object] {
	% 	\begin{tabular}{p{\nodeTextWidthRatio\textwidth}} 
\className{Wheel}

\desc{note that we could but don’t have to  model a wheel much. Just an object which evolves in its position}
\\ \hline
\data{angularSpeed}
\\ \hline
\method{evolve(delta-t)}
\end{tabular} 


	% };

	% \node[exampleClass ] (DiningTable) [right=of Wheel] {
	% 	\begin{tabular}{p{\nodeTextWidthRatio\textwidth}} 
\className{DiningTable}
\\ \hline
\data{Cylinder legs[]}

\data{Cube topSurface}
\\ \hline
\end{tabular} 


	% };

	\node[exampleClass ] (VaccumCleaner) [below=of Robot] {
		\begin{tabular}{p{\nodeTextWidthRatio\textwidth}} 
\className{VaccumCleaner}

\desc{in version 0, a cylinder with no wheels and stuff.}
\\ \hline
\\ \hline
\end{tabular} 


	};

	\node[exampleClass ] (Box) [below=of Object] {
		\begin{tabular}{p{\nodeTextWidthRatio\textwidth}} 
\className{Box}
\\ \hline
\\ \hline
\end{tabular} 


	};

	\node[exampleClass ] (BumperSensor) [below=of Sensor] {
		\begin{tabular}{p{\nodeTextWidthRatio\textwidth}} 
\className{BumperSensor}

\desc{detects at a hit upon intersection, boolean result}
\\ \hline
\\ \hline
\end{tabular} 


	};

	% \node[exampleClass ] (LidarSensor) [left=of BumperSensor] {
	% 	\begin{tabular}{p{\nodeTextWidthRatio\textwidth}} 
\className{LidarSensor}

\desc{essentially a light beam rotating at an x-y plane and measuring distance with surrounding objects}
\\ \hline
\\ \hline
\end{tabular} 


	% };

	\node[exampleClass, wideBox ] (RigidPointBall) [right=of RigidPhysicalObject] {
		\begin{tabular}{p{\nodeTextWidthRatio\textwidth}} 
\className{RigidPointBall}

\desc{A ball with relevant calculations for reflection upon bump, but not rotating}

\desc{v0: 2D cylinder instead of sphere.}
\\ \hline
\\ \hline
\end{tabular} 
    
    
	};

	% \draw [inheritedFrom] (Disk) -- (Shape);
	\draw [inheritedFrom] (Cylinder) -- (Shape);
	\draw [inheritedFrom] (Cube) -- (Shape);
	% \draw [inheritedFrom] (DiningTable) -- (Object);
	\draw [inheritedFrom] (Box) -- (Object);
	\draw [inheritedFrom] (VaccumCleaner) -- (Robot);
	\draw [inheritedFrom] (BumperSensor) -- (Sensor);
	% \draw [inheritedFrom] (LidarSensor) -- (Sensor);
	\draw [inheritedFrom] (RigidPointBall) -- (RigidPhysicalObject);

\fi

%%%%% the matrix annotation part :D

%\node[inner sep=0pt, right of= InIn, xshift=-6cm, yshift=1cm] (neo)
%    {\includegraphics[width=.25\textwidth]{neo.jpg}};

%\node[inner sep=0pt, right of= neo, xshift=-4.5cm] (architect)
%    {\includegraphics[width=.25\textwidth]{architect.jpg}};

%\node [rectangle,black,draw, below of=architect, yshift=-2cm, text width = 10cm] 
%    {\textbf{Neo} : Why am I here?\\
%	\textbf{The Architect} : Your life is the sum of a remainder of an unbalanced equation inherent to the programming of the matrix. You are the eventuality of an anomaly, which despite my sincerest efforts I have been unable to eliminate from what is otherwise a harmony of mathematical precision.};

% \node[inner sep=0pt, right of= Map, xshift=5cm] (smith)
%     {\includegraphics[width=.4\textwidth]{smith.jpg}};

% \node [rectangle,black,draw, below of=smith, yshift=-1.5cm, text width = 8cm] 
%     {\textbf{Smith to Neo}: Because Of You, I'm No Longer An Agent Of This System. Because Of You, I've Changed. I'm Unplugged. A New Man, So To Speak.};


\end{tikzpicture}

\end{document}
